\documentclass[12pt]{article}
%---DOCUMENT MARGINS---
\usepackage{geometry} % Required for adjusting page dimensions and margins
\geometry{
	paper=a4paper, % Paper size, change to letterpaper for US letter size
	top=4cm, % Top margin
	bottom=2cm, % Bottom margin
	left=2cm, % Left margin
	right=2cm, % Right margin
	headheight=3cm, % Header height
	%footskip=1.5cm, % Space from the bottom margin to the baseline of the footer
	headsep=0.5cm, % Space from the top margin to the baseline of the header
	%showframe, % Uncomment to show how the type block is set on the page
}
\usepackage{cmap}

\usepackage[T1,T2A]{fontenc}
\usepackage{fontspec}
\setmainfont[
	BoldFont={* Bold},
	ItalicFont={* Italic},
	BoldItalicFont={* BoldItalic}
]{DejaVu Serif Condensed}
\setsansfont{DejaVu Sans}
\setmonofont{Dejavu Sans Mono}
\usepackage[math-style=TeX]{unicode-math}
\setmathfont{Latin Modern Math}

\usepackage[nobottomtitles*]{titlesec}
\titleformat
{\section} % command
{\fontsize{14}{14}\sffamily\bfseries} % format
{} % label
{0pt} % sep
{} % before-code
\titlespacing{\section}{0pt}{0.75em}{0.25em}
\newcommand{\tl}{}
\newcommand{\ml}{}
\newcommand{\problem}[1]{%
	\section[#1]{#1 \fontsize{14}{14}\selectfont{\hfill{\normalfont{
				\emoji{hourglass-not-done}\tl \space\space \emoji{floppy-disk} \ml}}}}
}
\AddToHook{cmd/section/before}{\clearpage}

\usepackage[dvipsnames]{xcolor}
\titleformat
{\subsection} % command
{\fontsize{14}{14}\sffamily\bfseries} % format
{\includesvg[width=0.035\textwidth, inkscapelatex=false]{lion}\space} % label
{0pt} % sep
{} % before-code
\titlespacing{\subsection}{0pt}{0.75em}{0.25em}

\setlength{\parskip}{0.5em}
\setlength{\parindent}{0pt}
\renewcommand{\baselinestretch}{1.2}
\sloppy

\usepackage{listings}
\definecolor{lstbg}{RGB}{250,250,252}
\definecolor{lstframe}{RGB}{220,220,225}
\lstdefinestyle{mystyle}{
	backgroundcolor=\color{lstbg},
	frame=single,
	rulecolor=\color{lstframe},
	basicstyle=\ttfamily\small,
	keywordstyle=\bfseries,
	breaklines=true,
	aboveskip=0.5\baselineskip,
	belowskip=0\baselineskip  
}
\lstset{style=mystyle, language=C++}

\usepackage{fancyhdr}
\pagestyle{fancy}
\setlength{\headwidth}{\textwidth}
\newcommand{\round}{}
\newcommand{\problemName}{}
\newcommand{\lang}{English}
\fancyhead[L]{
	\begin{minipage}{0.4\textwidth}
		\bf{
			EJOI 2025 \round\\
			\problemName\\
			\emoji{globe-with-meridians} \lang
		}
	\end{minipage}
	\vspace{0.5cm}
}
\fancyhead[R]{%
	\begin{minipage}{0.6\textwidth}
		\includesvg[width=\textwidth, inkscapelatex=false]{logo}
	\end{minipage}
	\vspace{0.5cm}
}
\usepackage{lastpage}
\fancyfoot[C]{\problemName\space(\lang) \thepage\ / \pageref*{LastPage}}

\raggedbottom

\usepackage{amsmath}
\usepackage{stmaryrd}

\usepackage{graphicx}
\graphicspath{{./}}
\usepackage[export]{adjustbox}
\usepackage{wrapfig}
\makeatletter
\patchcmd\WF@putfigmaybe{\lower\intextsep}{}{}{\fail}
\AddToHook{env/wrapfigure/begin}{\setlength{\intextsep}{0pt}}
\makeatother
\usepackage[inkscapearea=page, inkscapepath=./svg-inkscape]{svg}
\svgpath{{./}}

\usepackage{makecell}
\usepackage{tabularray}
\AtBeginEnvironment{table}{\vspace{-0.2cm}}
\AtEndEnvironment{table}{\vspace{-0.2cm}}
\usepackage{float}
\usepackage{placeins}
\usepackage{caption}
\captionsetup[table]{
	skip=0.25em,
	singlelinecheck=false, justification=justified
}

\usepackage{enumitem}
\setlist{itemsep=-0.4em, leftmargin=22pt, topsep=-\parskip}
\AfterEndEnvironment{itemize}{\vspace{\parskip}}
\newcommand{\tabitem}{\indent~~\llap{\textbullet}~~}

\usepackage{hyperref}
\hypersetup{
	colorlinks=true,
	citecolor=blue,
	linkcolor=blue,
	urlcolor=cyan,
}

\usepackage{emoji}

\usepackage[ukrainian]{babel}

\begin{document}

\renewcommand{\round}{Day 1}
\renewcommand{\problemName}{Problem Prison}
\renewcommand{\lang}{Ukrainian}

\renewcommand{\tl}{$2$ sec.}
\renewcommand{\ml}{$1024$ MB}
\problem{Задача В'язниця}

Алісу та Боба несправедливо засудили до в'язниці суворого режиму. Тепер вони повинні спланувати свою втечу. Для цього їм потрібно мати можливість спілкуватися якомога ефективніше (зокрема, Алісі потрібно щодня надсилати Бобу інформацію). Однак вони не можуть зустрічатися і можуть обмінюватися інформацією лише через записки, написані на серветках. Щодня Аліса хоче надсилати Бобу новий фрагмент інформації – число від $0$ до $N-1$. На кожному обіді Аліса отримує три серветки та пише на кожній серветці число від $0$ до $M-1$ (можливі повторення) та залишає їх на своєму місці. Потім їхній ворог, Чарлі, знищує одну із серветок і змішує дві інші. Зрештою, Боб знаходить дві серветки, що залишилися, і зчитує числа на них. Він повинен точно розшифрувати початкове число, яке Аліса хотіла йому надіслати. На серветках обмежений простір, тому $M$ фіксоване. Однак мета Аліси та Боба – максимізувати пропускну здатність інформації, тому вони можуть вільно вибирати $N$ якомога більшим. Допоможіть Алісі та Бобу, застосувавши для кожного з них стратегію, намагаючись максимізувати значення $N$.

\subsection{Деталі реалізації}
Оскільки це проблема зв'язку, ваша програма буде виконуватися у двох окремих виконаннях (одне для Аліси та одне для Боба), які не можуть обмінюватися даними або спілкуватися будь-яким способом, окрім описаного тут. Вам потрібно реалізувати три функції:

\begin{lstlisting}
 int setup(int M);
\end{lstlisting}

Вона буде викликано один раз на початку виконання вашої програми Алісою та один раз на початку виконання Бобом. Їй надається $M$, і вона повинна повернути бажане $N$. Обидва виклики \texttt{setup} повинні повернути одне і те саме $N$.

\begin{lstlisting}
 std::vector<int> encode(int A);
\end{lstlisting}

Вона реалізує стратегію Аліси. Вона буде викликана з числом для кодування $A$ ($0 \leq A < N$) і повинна повернути три числа $W_1, W_2, W_3$ ($0 \leq W_i < M$), які кодують $A$. Ця функція буде викликана загалом $T$ разів -- один раз на день (значення $A$ можуть повторюватися між днями).

\begin{lstlisting}
 int decode(int X, int Y);
\end{lstlisting}

Вона реалізує стратегію Боба. Вона буде викликана з двома з трьох чисел, повернутих \texttt{encode}, у певному порядку. Вона повинна повернути те саме значення $A$, яке отримала \texttt{encode}. Ця функція також буде викликана $T$ разів -- відповідно до $T$ викликів \texttt{encode}; вони будуть у тому ж порядку. Усі виклики \texttt{encode} відбуватимуться перед усіма викликами \texttt{decode}.

\subsection{Обмеження}

\begin{itemize}
\item $M \leq 4300$
\item $T = 5000$
\end{itemize}

\subsection{Підрахунок балів}

Для конкретного підзавдання частка S отриманих вами балів залежить від найменшого значення $N$, яке повертає \texttt{setup} у будь-якому тесті в цьому підзавданні. Вона також залежить від $N^*$, яке є цільовим значенням $N$, необхідним для отримання повної кількості балів за підзадачу:

\begin{itemize}
\item Якщо ваше рішення не пройде хоча б один тест, тоді $S = 0$.
\item Якщо $N \geq N^*$, тоді $S = 1.0$.
\item Якщо $N < N^*$, тоді $S = \max\left(0.35 \max\left(\frac{\log(N) - 0.985 \log(M)}{\log(N^*) - 0.985 \log(M)}, 0.0\right)^{0.3} + 0.65 \left(\frac{N}{N^*}\right)^{2.4}, 0.01\right)$.
\end{itemize}

\subsection{Підзадачі}

\begin{table}[H]
\begin{tblr}{|Q[c,m]|Q[c,m]|Q[c,m]|Q[c,m]|}
\hline
\textbf{Підзадача} & \textbf{Бали} & $M$ & $N^*$ \\
\hline
$1$ & $10$ & $700$ & $82017$ \\
\hline
$2$ & $10$ & $1100$ & $202217$ \\
\hline
$3$ & $10$ & $1500$ & $375751$ \\
\hline
$4$ & $10$ & $1900$ & $602617$ \\
\hline
$5$ & $10$ & $2300$ & $882817$ \\
\hline
$6$ & $10$ & $2700$ & $1216351$ \\
\hline
$7$ & $10$ & $3100$ & $1603217$ \\
\hline
$8$ & $10$ & $3500$ & $2043417$ \\
\hline
$9$ & $10$ & $3900$ & $2536951$ \\
\hline
$10$ & $10$ & $4300$ & $3083817$ \\
\hline
\end{tblr}
\end{table}
\FloatBarrier

\subsection{Приклад}

Розглянемо наступний приклад з $T = 5$. Тут ми маємо схему кодування, де Аліса надсилає три однакові числа для кодування $0$ або три різні числа для кодування $1$. Зверніть увагу, що Боб може декодувати вихідне число з будь-яких двох із трьох чисел, надісланих Алісою.

\begin{table}[H]
\begin{tblr}{|Q[c,m]|Q[c,m]|Q[c,m]|}
\hline
\textbf{\makecell{Виконання}} & \textbf{\makecell{Виклик функції}} & \textbf{Повернене значення} \\
\hline
Аліса & \texttt{setup(10)} & \texttt{2} \\
\hline
Боб & \texttt{setup(10)} & \texttt{2} \\
\hline
Аліса & \texttt{encode(0)} & \texttt{\{5, 5, 5\}} \\
\hline
Аліса & \texttt{encode(1)} & \texttt{\{8, 3, 7\}} \\
\hline
Аліса & \texttt{encode(1)} & \texttt{\{0, 3, 1\}} \\
\hline
Аліса & \texttt{encode(0)} & \texttt{\{7, 7, 7\}} \\
\hline
Аліса & \texttt{encode(1)} & \texttt{\{6, 2, 0\}} \\
\hline
Боб & \texttt{decode(5, 5)} & \texttt{0} \\
\hline
Боб & \texttt{decode(8, 7)} & \texttt{1} \\
\hline
Боб & \texttt{decode(3, 0)} & \texttt{1} \\
\hline
Боб & \texttt{decode(7, 7)} & \texttt{0} \\
\hline
Боб & \texttt{decode(2, 0)} & \texttt{1} \\
\hline
\end{tblr}
\end{table}

\subsection{Приклад градера}

Для зразка оцінювання всі виклики \texttt{encode} та \texttt{decode} будуть виконуватися в одному і тому ж виконанні вашої програми. Крім того, \texttt{setup} буде викликатися тільки один раз (на відміну від двох разів, по одному разу на кожне виконання, як в системі оцінювання).

Вхідні дані складаються з одного цілого числа $M$. Потім він виведе число $N$, яке повернув ваш \texttt{setup}.
Потім він викличе \texttt{encode} та \texttt{decode} $T$ з випадково згенерованими числами від $0$ до $N-1$. Серед трьох чисел, які поверне \texttt{encode} він вибере 2 випадковим чином, в випадковому порядку і передасть їх в \texttt{decode}. Якщо ваше рішення спрацювало некоректно то він виведе повідомлення про помилку.

\end{document}